\begin{thebibliography}{25}
\interlinepenalty=10000

\bibitem{vitaldb}
H.-C. Lee, Y. Park, S.-B. Yoon, \emph{et al.}, ``VitalDB, a high-fidelity multi-parameter vital signs database in surgical patients,'' \emph{Scientific Data}, 9: 279, 2022.

\bibitem{mover}
M. Samad, M. Angel, J. Rinehart, \emph{et al.}, ``Medical Informatics Operating Room Vitals and Events Repository (MOVER): a public-access operating room database,'' \emph{JAMIA Open}, 6(4): ooad084, 2023.

\bibitem{inspire}
L. Lim, H. Lee, C.-W. Jung, \emph{et al.}, ``INSPIRE, a publicly available research dataset for perioperative medicine,'' \emph{Scientific Data}, 11: 655, 2024.

\bibitem{ricu}
N. Bennett, D. Ple\v{c}ko, I.-F. Ukor, \emph{et al.}, ``ricu: R's interface to intensive care data,'' \emph{GigaScience}, 12: giad041, 2023.

\bibitem{yaib}
R. van de Water, H. Schmidt, P. Elbers, \emph{et al.}, ``Yet Another ICU Benchmark: a flexible multi-center framework for clinical ML,'' in \emph{Proceedings of the 12th International Conference on Learning Representations (ICLR)}. Vienna, Austria: OpenReview, 2024.

\bibitem{ehrshot}
M. Wornow, R. Thapa, E. Steinberg, \emph{et al.}, ``EHRSHOT: an EHR benchmark for few-shot evaluation of foundation models,'' in \emph{Advances in Neural Information Processing Systems}, 36: 67125--67137. New Orleans, LA, USA: Curran Associates, Inc., 2023.

\bibitem{ohdsi}
G. Hripcsak, J. D. Duke, N. H. Shah, \emph{et al.}, ``Observational Health Data Sciences and Informatics (OHDSI): opportunities for observational researchers,'' in \emph{MEDINFO 2015: eHealth-enabled Health, Proceedings of the 15th World Congress on Health and Biomedical Informatics}, 216: 574--578. S\~ao Paulo, Brazil: IOS Press, 2015, doi: 10.3233/978-1-61499-564-7-574.

\bibitem{snomed}
SNOMED International, ``What is SNOMED CT?,'' website. [Online]. Available: \url{https://www.snomed.org/what-is-snomed-ct}. Accessed: Aug. 22, 2026.

\bibitem{umls}
O. Bodenreider, ``The Unified Medical Language System (UMLS): integrating biomedical terminology,'' \emph{Nucleic Acids Research}, 32(suppl. 1): D267--D270, 2004.

\bibitem{kim_omop_perioperative}
H. Kim, H. Lee, D. Lee, \emph{et al.}, ``Harmonizing perioperative open datasets using the OMOP-CDM via LLM and LangChain,'' conference presentation at the 10th OHDSI Global Symposium, New Brunswick, NJ, USA, Oct. 22--24, 2024. [Online]. Available: \url{https://www.ohdsi.org/wp-content/uploads/2024/12/8b.-LLM-and-OHDSI-Part-2_Hyeonsik-Kim.pdf}. Accessed: Aug. 22, 2026.

\bibitem{dnorm}
R. Leaman, R. Islamaj Do\u{g}an, and Z. Lu, ``DNorm: disease name normalization with pairwise learning to rank,'' \emph{Bioinformatics}, 29(22): 2909--2917, 2013.

\bibitem{bert_normalization}
Z. Ji, Q. Wei, and H. Xu, ``BERT-based ranking for biomedical entity normalization,'' in \emph{AMIA Joint Summits on Translational Science Proceedings}, pp. 269--277. Bethesda, MD, USA: American Medical Informatics Association, 2020.

\bibitem{n2c2}
Y.-F. Luo, S. Henry, Y. Wang, \emph{et al.}, ``The 2019 n2c2/UMass Lowell shared task on clinical concept normalization,'' \emph{Journal of the American Medical Informatics Association}, 27(10): 1529--e1, 2020, doi: 10.1093/jamia/ocaa106.

\bibitem{clinicalbert}
E. Alsentzer, J. R. Murphy, W. Boag, \emph{et al.}, ``Publicly available clinical BERT embeddings,'' in \emph{Proceedings of the 2nd Clinical Natural Language Processing Workshop}, pp. 72--78. Minneapolis, MN, USA: Association for Computational Linguistics, 2019.

\bibitem{rag}
P. Lewis, E. Perez, A. Piktus, \emph{et al.}, ``Retrieval-augmented generation for knowledge-intensive NLP tasks,'' in \emph{Advances in Neural Information Processing Systems}, 33: 9459--9474. Red Hook, NY, USA: Curran Associates, Inc., 2020.

\bibitem{llm_snomed}
S. Huh, ``Comparative analysis of ChatGPT-4 for automated mapping of local medical terminologies to SNOMED CT,'' in \emph{Intelligent Health Systems---From Technology to Data and Knowledge: Proceedings of MIE 2025}, 327: 813--817. Glasgow, UK: IOS Press, 2025, doi: 10.3233/SHTI250472.

\bibitem{faiss}
J. Johnson, M. Douze, and H. J\'egou, ``Billion-scale similarity search with GPUs,'' \emph{IEEE Transactions on Big Data}, 7(3): 535--547, 2021.

\bibitem{minibatch_kmeans}
D. Sculley, ``Web-scale K-means clustering,'' in \emph{Proceedings of the 19th International Conference on World Wide Web}, pp. 1177--1178. Raleigh, NC, USA: Association for Computing Machinery, 2010.

\bibitem{hdbscan}
R. J. G. B. Campello, D. Moulavi, and J. Sander, ``Density-based clustering based on hierarchical density estimates,'' in \emph{Advances in Knowledge Discovery and Data Mining: PAKDD 2013, Part II}, \emph{Lecture Notes in Computer Science}, 7819: 160--172. Gold Coast, QLD, Australia: Springer Berlin Heidelberg, 2013.

\bibitem{cox_logistic}
D. R. Cox, ``The regression analysis of binary sequences,'' \emph{Journal of the Royal Statistical Society: Series B (Methodological)}, 20(2): 215--232, 1958.

\bibitem{breiman_rf}
L. Breiman, ``Random forests,'' \emph{Machine Learning}, 45(1): 5--32, 2001.

\bibitem{xgboost}
T. Chen and C. Guestrin, ``XGBoost: a scalable tree boosting system,'' in \emph{Proceedings of the 22nd ACM SIGKDD International Conference on Knowledge Discovery and Data Mining}, pp. 785--794. San Francisco, CA, USA: Association for Computing Machinery, 2016.

\bibitem{backprop}
D. E. Rumelhart, G. E. Hinton, and R. J. Williams, ``Learning representations by back-propagating errors,'' \emph{Nature}, 323: 533--536, 1986.

\bibitem{lecun_cnn}
Y. LeCun, L. Bottou, Y. Bengio, \emph{et al.}, ``Gradient-based learning applied to document recognition,'' \emph{Proceedings of the IEEE}, 86(11): 2278--2324, 1998.

\bibitem{tabtransformer}
X. Huang, A. Khetan, M. Cvitkovic, \emph{et al.}, ``TabTransformer: tabular data modeling using contextual embeddings,'' arXiv preprint arXiv:2012.06678, 2020.

\end{thebibliography}
